\chapter{Preliminaries on Special Embedded Matrix Lie Groups}
\label{chapter:preliminaries_lie_groups}

This chapter summarizes the Lie group notation required for the continuous 
and discrete rigid-body mechanics developed in Chapters \ref{chapter:continous_mechanics_of_rigid_bodies_on_lie_groups} and \ref{chapter:discrete_mechanics_for_rigid_bodies_on_lie_groups}, following the definitions of Lie groups and Lie algebras in \cite{Hall2000LieGroups} and the geometric-mechanics notation in \cite{Lee2008}.

\section{Special Matrix Lie Groups}
\label{sec:special_matrix_lie_groups}

A Lie group is a smooth manifold that is also a group, where multiplication and inversion are smooth maps. A real matrix Lie group is a closed subgroup
$G\subset\mathrm{GL}(n,\mathbb{R})$. Its group operation is matrix
multiplication, and it is an embedded smooth manifold for which
multiplication and inversion are smooth maps
\cite{Hall2000LieGroups}.

Rotations in \(n\) dimensions are represented by the special orthogonal group
\begin{equation}
    \mathrm{SO}(n)
    =
    \left\{
        \bm{R}\in\mathbb{R}^{n\times n}
        \;\middle|\;
        \bm{R}^{\top}\bm{R}=\bm{I}_{n},\;
        \det(\bm{R})=1
    \right\}.
    \label{eq:so_group}
\end{equation}
The groups \(\mathrm{SO}(2)\) and \(\mathrm{SO}(3)\) describe planar and spatial rotations \cite{Lee2008}.

Combined rotational and translational motion is represented by the special Euclidean group
\begin{equation}
    \mathrm{SE}(n)
    =
    \left\{
        \bm{T}
        =
        \begin{bmatrix}
            \bm{R} & \bm{p}\\
            \bm{0} & 1
        \end{bmatrix}
        \;\middle|\;
        \bm{R}\in\mathrm{SO}(n),\;
        \bm{p}\in\mathbb{R}^{n}
    \right\}.
    \label{eq:se_group}
\end{equation}
Here, \(\bm{R}\) denotes the attitude and \(\bm{p}\) the position \cite{Hall2000LieGroups}. However, in this thesis 
rotations and translations are represented separately
in maximal coordinates, elaborated in Section \ref{sec:multi_body_systems}. Therefore, the derivations mainly use
$\mathrm{SO}(n)$ and its Lie algebra $\mathrm{so}(n)$.

\section{Special Lie Algebras}
\label{sec:special_lie_algebras}

\paragraph{Lie algebras and the exponential map}
\cite{Hall2000LieGroups} defines the Lie algebra \(\mathfrak{g}\) of a matrix Lie group \(G\) as its tangent space at the identity,
\begin{equation}
    \mathfrak{g}=T_{e}G.
    \label{eq:lie_algebra_tangent_space}
\end{equation}
For matrix Lie groups, the Lie bracket $[\cdot,\cdot]: \mathfrak{g} \times \mathfrak{g} \to \mathfrak{g}$ 
is defined as
\begin{equation}
    [\bm{a},\bm{b}]
    =
    \bm{a}\bm{b}-\bm{b}\bm{a},
    \qquad
    \bm{a},\bm{b}\in\mathfrak{g},
    \label{eq:matrix_lie_bracket}
\end{equation}
which is bilinear and skew-symmetric \cite{Hall2000LieGroups}. The Lie algebras corresponding to \(\mathrm{SO}(n)\) and \(\mathrm{SE}(n)\) are
\begin{align}
    \mathfrak{so}(n)
    &=
    \left\{
        \bm{\omega}\in\mathbb{R}^{n\times n}
        \;\middle|\;
        \bm{\omega}^{\top}=-\bm{\omega}
    \right\}
    \label{eq:so_algebra},
    \qquad
    &\mathfrak{so}(n) = T_{e}\mathrm{SO}(n)
    \\
    \mathfrak{se}(n)
    &=
    \left\{
        \begin{bmatrix}
            \bm{\omega} & \bm{v}\\
            \bm{0}^{\top} & 0
        \end{bmatrix}
        \;\middle|\;
        \bm{\omega}\in\mathfrak{so}(n),\;
        \bm{v}\in\mathbb{R}^{n}
    \right\},
    \qquad
    &\mathfrak{se}(n) = T_{e}\mathrm{SE}(n).
    \label{eq:se_algebra}
\end{align}
For \(\mathrm{SO}(2)\), \cite{Lee2008} defines the hat map $(\cdot)^\wedge: \mathbb{R} \rightarrow \mathfrak{so}(2)$ as
\begin{equation}
    \widehat{\omega}
    =
    \begin{bmatrix}
        0 & -\omega\\
        \omega & 0
    \end{bmatrix},
    \label{eq:hat_map_so2}
\end{equation}
with $\omega \in \mathbb{R}$.
Additionally, for \(\mathrm{SO}(3)\), the hat map \((\cdot)^{\wedge}:\mathbb{R}^{3}\rightarrow\mathfrak{so}(3)\) is defined by
\begin{equation}
    \widehat{\bm{\omega}}
    =
    \begin{bmatrix}
        0 & -\omega_{3} & \omega_{2}\\
        \omega_{3} & 0 & -\omega_{1}\\
        -\omega_{2} & \omega_{1} & 0
    \end{bmatrix},
    \qquad
    \widehat{\bm{\omega}}\bm{b}=\bm{\omega}\times\bm{b},
    \label{eq:hat_map_so3}
\end{equation}
with $\bm{\omega}, \bm{b} \in \mathbb{R}^3$ and the corresponding inverse map \((\cdot)^{\vee}: \mathfrak{so}(3) \rightarrow \mathbb{R}^3\).

The matrix exponential of an $n\times n$ matrix $\bm{X}$
\begin{equation}
    \exp(\bm{X})=\sum_{m=0}^{\infty} \frac{\bm{X}^m}{m!}
    \label{eq:matrix_exponential}
\end{equation}
defines the exponential map \(\exp:\mathfrak{g}\rightarrow G\) as a local diffeomorphism between a neighborhood of \(\bm{0}\in\mathfrak{g}\) and a neighborhood of \(e\in G\) \cite{Hall2000LieGroups,Lee2008}.

\paragraph{Left and right translations}
For \(\bm{g},\bm{h}\in G\), the left and right translation maps are
\begin{equation}
    \mathbb{L}_{\bm{h}}(\bm{g})=\bm{h}\bm{g},
    \qquad
    R_{\bm{h}}(\bm{g})=\bm{g}\bm{h}.
    \label{eq:left_right_translation}
\end{equation}
Thereby, \cite{Lee2008} states that for a differentiable matrix-valued curve $\bm{g}(t) \in G$, the left-trivialized or body-fixed velocity $\bm{\xi} \in \mathfrak{g}$ is given by
\begin{equation}
    \bm{\xi}=\bm{g}^{-1}\dot{\bm{g}},
    \qquad
    \dot{\bm{g}}=\bm{g}\bm{\xi}.
    \label{eq:left_trivialized_velocity}
\end{equation}

\begin{proposition}[Differential of left translation on a matrix Lie group: $T_{\bm{x}}\mathbb{L}_{\bm{h}}$]
Let $G\subset GL(n,\mathbb{R})$ be a matrix Lie group.
For $\bm{h}\in G$, define $\mathbb{L}_{\bm{h}}:G\to G$ by $\mathbb{L}_{\bm{h}}(\bm{g})=\bm{h}\bm{g}$ \cite{Hall2000LieGroups}.
Then for each $\bm{x}\in G$ the differential
\[
T_{\bm{x}}\mathbb{L}_{\bm{h}}: T_{\bm{x}}G \to T_{\bm{h}\bm{x}}G,
\]
of the left translation map $\mathbb{L}_{\bm{h}}:G\to G$ is given by
\begin{align}\label{eq:preliminaries_differential_left_translation}
T_{\bm{x}}\mathbb{L}_{\bm{h}}(\bm{v})=\bm{h}\bm{v}.
\end{align}
\end{proposition}
\begin{proof}
By the definition of the tangent space of an embedded submanifold given in \cite{Hall2000LieGroups}, there exists a smooth matrix-valued curve $\bm{\gamma}: (-\epsilon,\epsilon) \to G$ with $\bm{\gamma}(0)=\bm{x}$ and the tangent matrix $\dot{\bm{\gamma}}(0)=\bm{v} \in T_{\bm{x}}G$.
In \cite{Lee2008}, the differential is defined as
\[
T_{\bm{x}}\mathbb{L}_{\bm{h}}(\bm{v})=\left.\frac{d}{dt}\mathbb{L}_{\bm{h}}(\bm{\gamma}(t))\right|_{t=0}
=\left.\frac{d}{dt}(\bm{h}\bm{\gamma}(t))\right|_{t=0}=\bm{h}\dot{\bm{\gamma}}(0),
\]
with the left translation map $\mathbb{L}_{\bm{h}}(\bm{g})=\bm{h}\bm{g}$. Since $\bm{h}$ is constant and $\dot{\bm{\gamma}}(0)=\bm{v}$, it follows that
\[
T_{\bm{x}}\mathbb{L}_{\bm{h}}(\bm{v})=\bm{h}\dot{\bm{\gamma}}(0)=\bm{h}\bm{v}.
\]
\end{proof}

It can be observed that $\mathbb{L}_{\bm{h}}(\bm{g})=\bm{h}\bm{g}$ and $T_{\bm{x}}\mathbb{L}_{\bm{h}}(\bm{v})=\bm{h}\bm{v}$ have the same matrix expression, but act on different objects. The left translation map $\mathbb{L}_{\bm{h}}(\bm{g})$ multiplies group elements $\bm{g} \in G$, while its differential is the pushforward of tangent matrices $\bm{v} \in T_{\bm{x}}G \subset \mathbb{R}^{n \times n}$. The differential of a left translation can be applied to the special embedded matrix Lie groups $SO(n)$ and their corresponding Lie algebras $\mathfrak{so}(n)$. 

\paragraph{Differential of left translation $T_{\bm{R}}\mathbb{L}_{\bm{H}}$ on $SO(n)$}

The left translation map in $SO(n)$ is defined as $\mathbb{L}_{\bm{H}}(\bm{R}) = \bm{H}\bm{R}$ with $\bm{H},\bm{R} \in G$. Applied to eq. \eqref{eq:preliminaries_differential_left_translation} with $\bm{v} \in T_{\bm{R}}G$, it follows that
\begin{align*}
    T_{\bm{R}}\mathbb{L}_{\bm{H}}(\bm{v})= \bm{H}\bm{v}.
\end{align*}
The tangent matrix $\bm{v} \in T_{\bm{R}}G$ can be represented by $\bm{v} = T_{e}\mathbb{L}_{\bm{R}}(\bm{\omega}) = \bm{R}\bm{\omega}$ using the left translation map of $\bm{\omega} \in \mathfrak g := \mathfrak{so}(n)$, resulting in
\begin{align}\label{eq:preliminaries_differential_left_SO}
    T_{\bm{R}}\mathbb{L}_{\bm{H}}(\bm{R}\bm{\omega}) = \bm{H}(\bm{R}\bm{\omega}).
\end{align}

%\paragraph{Differential of left translation $T_g\mathbb{L}_h$ on $SE(n)$}

%The left translation map in $SE(n)$ is defined as $\mathbb{L}_{h}(g) = hg$ with $h,g \in G$. Applied to eq. \eqref{eq:preliminaries_differential_left_translation} with $v \in T_{g}G$, it follows that
%\begin{align*}
%    T_{g}\mathbb{L}_{h}(v)= hv.
%\end{align*}
%The tangent vector $v \in T_gG$ can be represented by $v = T_{e}\mathbb{L}_{g}(\xi) = g\xi$ using the left translation map of $\xi \in \mathfrak g := \mathfrak{se}(n)$ resulting in
%\begin{align}\label{eq:preliminaries_differential_left_SE}
%    T_g\mathbb{L}_h(g \xi) = h(g \xi).
%\end{align}



\paragraph{Adjoint and coadjoint operators}
\cite{Lee2008} introduces the adjoint operator $\operatorname{Ad}_{\bm{g}}: \mathfrak{g} \to \mathfrak{g}$ as
\begin{equation}
    \operatorname{Ad}_{\bm{g}}\bm{\eta}
    =
    \bm{g}\bm{\eta}\bm{g}^{-1}.
    \label{eq:adjoint_matrix}
\end{equation}
Its derivative  $\operatorname{ad}_{\bm{\xi}}: \mathfrak{g} \to \mathfrak{g}$ is defined as
\begin{equation}
    \operatorname{ad}_{\bm{\xi}}\bm{\eta}
    =
    \left.\frac{\mathrm{d}}{\mathrm{d}\epsilon}\right|_{\epsilon=0}
    \operatorname{Ad}_{\exp(\epsilon\bm{\xi})}\bm{\eta}
    =
    [\bm{\xi},\bm{\eta}]
    = \bm{\xi}\bm{\eta} - \bm{\eta}\bm{\xi},
    \label{eq:ad_lie_bracket}
\end{equation}
using the Lie bracket $[\cdot,\cdot]$ defined in \eqref{eq:matrix_lie_bracket}.

Let \(\mathfrak{g}^{*}\) be the dual space of \(\mathfrak{g}\). The notation $\langle\bm{\mu},\bm{\xi}\rangle = \bm{\mu}(\bm{\xi})$ denotes the pairing of a covector \(\bm{\mu}\in\mathfrak{g}^{*}\) with a vector \(\bm{\xi}\in\mathfrak{g}\) \cite{Lee2008}. 
Using the pairing, the coadjoint operator \(\operatorname{Ad}_{\bm{g}}^{*}:\mathfrak{g}^{*}\rightarrow\mathfrak{g}^{*}\) is defined as 
\begin{equation}
    \left\langle
        \operatorname{Ad}_{\bm{g}}^{*}\bm{\mu},
        \bm{\xi}
    \right\rangle
    =
    \left\langle
        \bm{\mu},
        \operatorname{Ad}_{\bm{g}}\bm{\xi}
    \right\rangle.
    \label{eq:coadjoint_definition}
\end{equation}
Its derivative \(\operatorname{ad}_{\bm{\eta}}^{*}:\mathfrak{g}^{*}\rightarrow\mathfrak{g}^{*}\) is given by
\begin{align}
    \left\langle\operatorname{ad}_{\bm{\eta}}^{*}\bm{\mu},\bm{\xi}\right\rangle
    &=
    \left\langle\bm{\mu},\operatorname{ad}_{\bm{\eta}}\bm{\xi}\right\rangle
    =
    \left\langle\bm{\mu},[\bm{\eta},\bm{\xi}]\right\rangle
    =
        \left\langle\bm{\mu},\bm{\eta}\bm{\xi} - \bm{\xi}\bm{\eta}\right\rangle
.
    \label{eq:coad_definition}
\end{align}

Additionally, the pairing can be used to define the inner product on the dual space of the Lie algebra \cite{Lee2008}. Using the trace inner product, $\mathfrak{so}(n)$ is identified with its
dual. For $n\in\{2,3\}$, the hat map gives
\begin{equation}
    \left\langle\widehat{\bm a},\widehat{\bm b}\right\rangle
    =
    \frac{1}{2}\operatorname{tr}\!\left(
        \widehat{\bm a}^{\top}\widehat{\bm b}
    \right)
    =
    \bm a^{\top}\bm b = \bm{a} \cdot \bm{b},
    \label{eq:trace_inner_product}
\end{equation}
where $\bm a,\bm b\in\mathbb{R}$ for $n=2$ and
$\bm a,\bm b\in\mathbb{R}^{3}$ for $n=3$.


\paragraph{Variations on Matrix Lie Groups}

For $g\in G$ and $\eta\in\mathfrak{g}=T_eG$, define the local variation
curve $\gamma:(-c,c)\rightarrow G$ by
\begin{equation}
    \gamma(\epsilon)
    =
    g\exp(\epsilon\eta),
    \qquad c>0,
    \label{eq:variation_curve_matrix_lie_group}
\end{equation}
where $\epsilon$ is the variation parameter. Since
$\exp(\epsilon\eta)\in G$, the curve $\gamma(\epsilon)$ remains in
$G$ and satisfies $\gamma(0)=g$ \cite{Lee2008}. The corresponding
infinitesimal variation is
\begin{equation}
    \delta g
    =
    \left.
    \frac{d}{d\epsilon}
    \right|_{\epsilon=0}
    \gamma(\epsilon)
    =
    \left.
    \frac{d}{d\epsilon}
    \right|_{\epsilon=0}
    g\exp(\epsilon\eta)
    =
    g\eta
    =
    T_e\mathbb{L}_g(\eta)
    \in T_gG.
    \label{eq:infinitesimal_variation_matrix_lie_group}
\end{equation}


%Explain brackets
%Inner product 
%-> Pairs
